
\section{Einleitung}

\subsection{Problemstellung}

Dieses Kapitel präzisiert die zentrale Herausforderung dieser Arbeit. Ein selbstbalancierendes Fahrrad stellt aufgrund seiner Bauart ein nichtlinear instabiles System dar. 
Ohne regelnden Eingriff kippt ein Fahrrad im Stand oder bei niedriger Geschwindigkeit. Erst durch die Kombination 
aus gyroskopischen Effekten der rotierenden Räder, dem Nachlauf der Lenkgeometrie und einem aktiven Regler 
lässt sich eine stabile Fahrt realisieren. Die vorliegenden Arbeiten verfolgen daher das Ziel, ein autonomaus 
fahrendes Fahrrad zu entwickeln, das sich sowohl selbst balancieren als auch eine vorgegebene Strecke verfolgen 
kann. Hierfür wurde eine hochrealistische Simulationsumgebung in Webots aufgebaut, die eine duale Reglerarchitektur 
abbildet. Ein schneller Balance‑Controller, implementiert in C, stabilisiert den Rollwinkel mit einer Abtastzeit 
von 2\,ms (500\,Hz) und gibt einen Lenkwinkel an den Antrieb weiter. Parallel dazu extrahiert ein Vision‑Controller 
in Python mit niedrigerer Frequenz (20\,Hz) Informationen aus Kamerabildern, plant eine Trajektorie und überträgt 
entsprechende Lenkwinkel‑ und Geschwindigkeitsbefehle. Die Herausforderung besteht darin, diese beiden Regelschleifen 
so zu koppeln, dass das Fahrzeug stabil bleibt und zugleich auf visuelle Reize reagiert. Darüber hinaus soll die 
simulierte Physik so realitätsnah sein, dass die Ergebnisse auf reale Fahrversuche übertragbar sind. Die Arbeit 
adressiert damit die Verbindung von Regelungs‑ und Computer‑Vision‑Methoden für ein komplexes, unteraktuiertes Fahrzeug.






\paragraph{Abgrenzung.}
Die Arbeit bleibt simulationsbasiert (Webots) und betrachtet nur eine monokulare Kamera ohne dynamische
Hindernisumfahrung; die Geschwindigkeit ist vorgegeben, die Kamera starr montiert. Hardware-Effekte
(z.\,B. Sensordrift, Antriebsnachgiebigkeit) sowie Multi-Sensor-Fusion sind nicht Gegenstand der Arbeit.










  
%--------------------------------------------------
% Stand der Technik
%--------------------------------------------------
\subsection{Stand der Technik}
Die Entwicklung selbstbalancierender Fahrräder ist seit einigen Jahren Gegenstand intensiver Forschung. 
Bereits 2005 analysierte K.J. Åström und Anders Lennartsson \cite{astrom2005} die Fahrdynamik 
von Fahrrädern und zeigten, dass ein geeignetes 
Lenkverhalten, insbesondere das Hineinlenken in die Fallrichtung, zur Stabilisierung genutzt werden kann. 
Dieses Prinzip, oft als "Steer into the Fall" (STF) bezeichnet, bildet die Grundlage vieler Balance-Regelungen. 
Allgemein sind Fahrräder ohne Fahrer inhärent instabil und verhalten sich ähnlich einem umgekehrten Pendel, 
das aktiv ausbalanciert werden muss. Aufbauend auf diesem Verständnis wurden verschiedene 
Methoden entwickelt, um ein fahrerloses Fahrrad im Gleichgewicht zu halten. Die mit Abstand einfachste Variante 
nutzt die Lenksteuerung, den durch gezieltes Verdrehen des Lenkers wird das Rad unter den Schwerpunkt manövriert, 
um Kippbewegungen auszugleichen. Zahlreiche Prototypen erreichen auf diese Weise dynamische Stabilisierung, 
erfordern jedoch eine gewisse Vorwärtsgeschwindigkeit. So demonstrierten beispielsweise Andreo, D. and Cerone schon 2010
einen nichtlinearen Regler, der mittels Vorderradlenkung ein autonomes Fahrrad stabilisieren konnte \cite{guo2010}. 
Ähnlich zeigte Deng et al. (2018), dass bereits ein einfacher PID- bzw. LQR-basierter Lenkwinkelregler 
\footnote{tEin LQR-Regler stabilisiert ein lineares System optimal, indem er Zustandsfehler und Stellaufwand minimiert.} ausreicht, 
um ein frei fahrendes Fahrrad aufrecht zu habitablen \cite{deng2018}. Andere Arbeiten nutzen erweiterte Verfahren: 
Andreo et al. (2009) realisierten eine Fahrstabilisierung mit einem LPV-Regler (linear parametrisch variierend), der 
verschiedene Fahrzustände abdecken konnte \cite{andreo2009}. Shafiekhani et al. (2015) implementierten einen neuro-fuzzy 
Regler, der durch adaptive Optimierung robuste Balance auch bei Modellunsicherheiten erzielte \cite{shafiekhani2015}. 
Darüber hinaus wurden Sliding-Mode-Verfahren (Gleitmodusregelung) vorgeschlagen, um gegen unbekannte Störungen unempfindlich zu sein
Insgesamt zeigt sich, dass klassische Regelungsansätze (PID, LQR) wie auch moderne Verfahren (MPC, adaptive und robuste Regelungen) 
erfolgreich zur Stabilisierung selbstbalancierender Fahrräder eingesetzt werden können \cite{kanjanawanishkul2015,shafiekhani2015,deng2018}.
Neben der Lenkung existieren alternative Mechanismen zur aktiven Stabilisierung. Historisch wurden vor allem zwei Ansätze untersucht
: Zum einen die Verwendung beweglicher Massen, zum anderen gyroskopische Effekte. Bewegte Masse: Hierbei wird ein Zusatzgewicht lateral 
verschoben oder als invertiertes Pendel im Rahmen angebracht, um durch Gewichtsverlagerung das Gleichgewicht zu halten. Diese Methode 
kann insbesondere bei sehr niedrigen Geschwindigkeiten oder im Stillstand für Stabilität sorgen, erhöht jedoch das Systemgewicht 
beträchtlich und ist dynamisch träge \cite{tamaldin2017}. Gyroskopische Stützung: Durch schnell rotierende Schwungmassen wird ein 
Kreiseleffekt erzeugt, der dem Umkippen entgegenwirkt. Zwei Varianten sind verbreitet: ein Reaktionsrad, das entlang der Längsachse 
montiert ist und durch Beschleunigen/Abbremsen Stabilisierungsmomente erzeugt, sowie ein Control Moment Gyroscope (CMG), bei dem 
ein rotierendes Schwungrad um eine Achse geschwenkt wird. Ein einzelnes Kreiselgerät erzeugt allerdings neben dem Stabilisierungs-Moment 
oft unerwünschte seitliche Drehmomente, weshalb Park und Yi (2020) beispielsweise ein gepaartes (scissored pair) Kreisel-System einsetzten, 
das Störmomente kompensiert \cite{park2020}. Mit einem solchen doppelten CMG konnte sogar stationäre Balance (ohne Vorwärtsfahrt) 
realisiert werden \cite{park2020}. Reaktionsräder wiederum sind in der Lage, ein Fahrrad aus dem Stand aufzurichten, erfordern 
dafür jedoch große Drehmomente und energiestarke Aktuatoren, was ihren Einsatzbereich begrenzt\cite{kanjanawanishkul2015}. 
Trotz dieser technischen Herausforderungen konnte Lam und Sin (2011) bereits ein kleines Roboterfahrrad mit Einzel-Schwungrad stabilisieren 
\cite{lam2011}. Tamaldin et al. (2017) präsentierten einen Gesamtentwurf eines selbstbalancierenden Fahrrads, der verschiedene solcher 
Stabilisierungsprinzipien in Betracht zog \cite{tamaldin2017}. Aktuell konzentriert sich die Forschung vermehrt darauf, Balancekontrolle 
mit Fahrtrajektorie zu verknüpfen – also ein selbstbalancierendes Fahrrad gezielt entlang einer vorgegebenen Spur oder Fahrspur zu steuern. 
Frühere Arbeiten adressierten bereits das Spurhalten: So gelang Saguchi et al. (2009) eine stabile Geradeausfahrt durch automatische 
Lenkregelung, während Aphiratsakun und Techakittiroj (2013) einen autonomen Fahrrad-Prototyp (AUSB2) vorstellten, der neben dem 
Balancieren auch einfache Trajektorien nachfahren konnte \cite{aphiratsakun2013}. Kanjanawanishkul (2015) kombinierte sogar eine 
stationäre Balance-Regelung mit einem MPC-basierten Pfadverfolger, sodass sein Fahrrad im Stillstand balancieren und bei Fahrt 
einer vorgegebenen Strecke folgen konnte \cite{kanjanawanishkul2015}. Ein weiterer Meilenstein ist die Arbeit von Xiong et al. (2024), 
die ein vollautonomes Fahrrad mit minimalem Aktuatoraufwand realisierten: Durch einen analytisch entworfenen linearen Lenkregler wurde 
das Fahrzeug sowohl im Geradeauslauf als auch im Kurvenfahren stabil gehalten, ohne dass zusätzliche Kreisel oder Gewichte nötig waren 
\cite{xiong2024}. Diese Ergebnisse zeigen, dass eine intelligente Lenkregelung prinzipiell ausreicht, um sowohl Balance als auch Spurhaltung 
zu gewährleisten – was für die angestrebte Fahrspurhaltung essentiell ist. Neben der Regelungstechnik spielt die Sensorik und Umfeldwahrnehmung 
eine wachsende Rolle im Stand der Technik. Für das autonome Fahren müssen Umgebung und Fahrspur erkannt werden. Hunziker (2019) koppelte z.B. 
Inertialsensoren mit GPS-Daten und setzte einen modellprädiktiven Regler ein, um ein elektrisches Fahrrad ohne Fahrer entlang eines vorgegebenen 
Kurses zu steuern \cite{Hunziker2019}. In neueren Projekten werden Kameras und Computer Vision integriert, um Fahrspuren oder Hindernisse zu 
detektieren. Ein prominentes Beispiel ist ein am Tsinghua-Universität entwickeltes autonomes Fahrrad, das mittels KI-Chip (Tianjic) sowohl 
selbstbalanciert als auch Objekterkennung und Sprachbefehle umsetzt. Wissenschaftlich dokumentiert ist die Arbeit von Ghidaei (2024), 
der ein video- und sensorgestütztes autonomes Fahrsystem für ein selbstbalancierendes E-Bike realisierte \cite{Ghidaei2024}. Dabei dienen 
Kamerabilder zur Spurhaltung (Lane Keeping), während zusätzliche Sensorfusion die Stabilität gewährleistet. Dieses Projekt stellt einen 
aktuellen Stand der Technik dar und zeigt, dass die Kombination aus Lageregelung (Balance-Controller) und Fahrspur-Assistent (Vision-basierte Steuerung) inzwischen machbar ist.
Abschließend ist festzuhalten, dass sich das Forschungsfeld rasant entwickelt. Zahlreiche Forschungsprototypen weltweit – 
vom LPV-geregelten Rad von Andreo (2009) bis zum neuesten visionbasierten E-Bike von Ghidaei (2024) – haben die Machbarkeit 
des selbstbalancierenden, spurgeführten Fahrrads demonstriert. Auch an der Goethe-Universität Frankfurt wurde in den letzten 
Jahren der Grundstein gelegt: Ranz (2021) entwickelte einen ersten selbstbalancierenden Fahrrad-Prototyp \cite{ranz2021bachelor}, 
Karabila (2022) implementierte eine Balance-Regelung dafür \cite{karabila2022}, und Zander (2023) verbesserte das Reglerdesign weiter 
\cite{zander2023}. Auf diesem Vorwissen aufbauend erweitert die vorliegende Arbeit den Stand der Technik um eine robuste Fahrspurhaltung, 
indem aktuelle Methoden der Regelungstechnik und Bildverarbeitung zusammengeführt werden.

%--------------------------------------------------
% Bisherige Arbeiten
%--------------------------------------------------
  \subsection{Bisherige Arbeiten}
  Die im vorherigen Abschnitt dargestellten Forschungsarbeiten zeigen, dass sowohl klassische 
  als auch moderne Regelungsverfahren prinzipiell geeignet sind, ein selbstbalancierendes 
  Fahrrad zu stabilisieren und mit Trajektorienfolgeregelungen zu kombinieren. Während sich 
  diese Studien überwiegend auf externe Forschungsgruppen beziehen, existieren auch an der 
  Goethe-Universität Frankfurt eine Reihe aufeinander aufbauender Abschlussarbeiten, die 
  wesentlich zur Entwicklung des hier untersuchten Systems beigetragen haben. 



  \subsubsection{Anna Ranz (2021) – Hardware-Grundlagen}

  Die Bachelorarbeit von Anna Ranz bildete den Startpunkt des Projekts. Ihr Ziel
  war die Konstruktion eines lauff\"ahigen Prototyps, mit dem grundlegende
  physikalische Eigenschaften eines selbstbalancierenden Fahrrads untersucht
  werden konnten. Kern der Hardware ist ein BeagleBone\,Black, der den
  Lenkmotor mittels MD49-Controller (EMG49) \"uber UART und die
  Inertialmess\-einheit BNO055 \"uber I\textsuperscript{2}C ansteuert. 
  Ergänzt wurde das System durch einen Nabenmotor am Hinterrad, der eine konstante 
  Vorwärtsfahrt sicherstellte. Damit standen erstmals sowohl die mechanische Plattform 
  als auch die notwendigen Sensor- und Aktor-Schnittstellen zur Verfügung, um 
  regelungstechnische Konzepte praktisch zu erproben.

  Ranz implementierte erste PID-Regelkreise, validierte die Sensorik in separaten
  Testprogrammen und analysierte die Dynamik des Gesamtsystems in Bezug auf
  Massentr\"agheit, Raddynamik und Lenkkinematik. Von besonderer Bedeutung war die 
  Untersuchung des Zusammenspiels zwischen Trägheitsmoment, Lenkwinkel und Rückstellkräften, 
  da diese Größen die Stabilität maßgeblich beeinflussen. Darüber hinaus wurden 
  Messreihen durchgeführt, um die Qualität der BNO055-IMU zu bewerten und deren Eignung 
  für die Rollwinkelbestimmung zu prüfen. Diese Kalibrierung stellte sicher, dass die 
  Daten mit ausreichend geringer Latenz und Genauigkeit für eine spätere Regelung 
  verwendet werden konnten.

  Die umfangreiche Dokumentation enth\"alt Schaltpl\"ane, 3D-Modelle und Messaufnahmen und
  bildet damit die hardwareseitige Grundlage f\"ur alle sp\"ateren Arbeiten 
  \cite{ranz2021bachelor}. Besonders hervorzuheben ist, dass Ranz nicht nur eine 
  funktionierende Testplattform aufbaute, sondern auch systematisch deren Grenzen 
  aufzeigte, etwa die eingeschränkte Reaktionsgeschwindigkeit des MD49-Motorcontrollers 
  oder die Latenz der UART-Kommunikation. Diese Erkenntnisse waren für die nachfolgenden 
  Arbeiten entscheidend, da sie Verbesserungsbedarf aufzeigten und konkrete Ansatzpunkte 
  für die Weiterentwicklung lieferten. 

  Damit markierte die Arbeit von Ranz den entscheidenden Übergang von theoretischen 
  Überlegungen hin zu einer realen Forschungsplattform: Erst durch die von ihr geschaffene 
  Hardware-Basis und die Analyse der physikalischen Dynamik war es möglich, in den 
  darauffolgenden Arbeiten sowohl Regelungsalgorithmen als auch softwareseitige 
  Optimierungen im praktischen Betrieb zu validieren.


  \subsubsection{Amine Karabila (2022) – Erste funktionierende Regelung}

  Amine Karabila schloss 2022 unmittelbar an den Hardware-Prototyp von Ranz an und 
  überführte das System von einzelnen Stellversuchen zu einem geschlossenen, 
  echtzeitfähigen Regelkreis. Kern­bestandteil ist eine dreistufige P-Regelung 
  auf dem BeagleBone Black, die den Rollwinkel der BNO055-IMU \cite{bno055} und 
  die Lenkerstellung des MD49-Encoders \cite{md49} auswertet, während ein hinterer 
  Nabenmotor die Vortriebsdrehzahl konstant hält \cite{karabila2022}. 
  Damit gelang erstmals eine softwareseitig realisierte Balancekontrolle, die den 
  Forschungsprototyp in eine echte Versuchsplattform verwandelte.
  
  Ein erstes Hindernis war die Strombegrenzung des MD49, die den EMG49-Lenkmotor 
  zu langsam ansprechen ließ. Karabila ersetzte deshalb den Leistungsteil durch eine 
  BTS7960-H-Brücke \cite{bts7960} und verwendete den MD49 nur noch zur hochauflösenden 
  Wegmessung. Nicht-blockierende UART-Aufrufe verkürzten die Encoder­abfrage von rund $4\,\mathrm{ms}$ 
  auf knapp $1\,\mathrm{ms}$, sodass die Gesamttransportzeit des Reglers auf etwa $2\,\mathrm{ms}$ sank. 
  Diese Optimierungen verdeutlichen die besonderen Herausforderungen eingebetteter Systeme: 
  Ein stabiler Regelkreis erfordert nicht nur gute Algorithmen, sondern auch eine 
  Hardware- und Softwarearchitektur, die minimale Verzögerungen gewährleistet.
  
  Die erste Regel­version koppelte die Motorgeschwindigkeit über einen gestreckten $\tanh$-Verlauf an 
  die Verweildauer des Rollwinkels in einer „habitablen Zone“. Außerhalb des Labors erwies sich 
  dieses Verfahren jedoch als zu träge: Der Lenker schlug regelmäßig an die Endanschläge, und Stützräder 
  mussten das Fahrrad vor dem Kippen bewahren. Die finale Lösung verknüpfte dagegen den gemessenen 
  Rollwinkel direkt mit einem geschwindigkeits­abhängigen Soll-Encoderwert. Dadurch pendelte der 
  Lenker nur noch mit etwa acht Ticks Amplitude um die Nullposition; ein integrierter Dämpfungsanteil 
  reduzierte die Motordrehzahl automatisch, sobald sich das System seiner Soll-Lage näherte. 
  
  Validiert wurde der Regler im Hof des Informatikgebäudes ($\approx 9.5\times26.5\,\text{m}$, fünf 
  Prozent Steigung). Bei Sollgeschwindigkeiten von $6$–$10\,\mathrm{km/h}$ blieb das Fahrrad während 
  aller aussagekräftigen Fahrten aufrecht; Tip-Over-Events traten nicht auf. Die erfolgreichste 
  Versuchsfahrt bei $10\,\mathrm{km/h}$ zeigte Spitzen von unter acht Encoder-Ticks und eine 
  Querabweichung, die optisch kaum wahrnehmbar war. Diese erfolgreiche Fahrt 
  dokumentiert die Funktionsfähigkeit der entwickelten Regelung. 
  
  Damit bewies Karabilas Arbeit erstmals die Realisierbarkeit einer selbst­balancierenden 
  Regelung und schuf die Grundlage für die PID-Optimierungen von Zander sowie die MPC-Navigation von Giray. 
  Insbesondere machte sie deutlich, dass die Regelgüte stark von niedrigen Latenzen und einer robusten 
  Kopplung zwischen Sensoren und Aktoren abhängt – ein zentrales Motiv, das auch in den nachfolgenden Arbeiten 
  beibehalten und weiterentwickelt wurde.
  
\newpage
\subsubsection{Jonah Zander (2023) – Optimierte PID-Balanceregelung}

Zander knüpfte 2023 an die Vorarbeiten an und formulierte als Kernziel, das Fahrrad mindestens 60\,s auf 
ebener Strecke ohne Bodenkontakt der Stützräder stabil fahren zu lassen \cite{zander2023}. Die Arbeit 
positioniert sich zwischen eingebetteter Echtzeit-Software und Regelungsentwurf: Latenz, Jitter und Sättigungseffekte 
sollten so adressiert werden, dass eine robuste Balance auch bei Richtungsänderungen möglich ist. 
Grundlage waren Ranz’ Hardware und Karabilas erste geschlossene Regelstrecke \cite{ranz2021bachelor, karabila2022}.

Zur Umsetzung entwickelte Zander eine modular aufgebaute C-Software auf dem BeagleBone\,Black (Debian), 
in der Multi-Threading Sensor-/Aktor-I/O, Reglerkern und Logging entkoppelt. Zeitkritische Abtast- und 
Vorverarbeitungsaufgaben wurden in die Programmable Realtime Units (PRUs) verlagert, wodurch deterministische 
Lesezyklen für Fahr- und Lenkwinkel im zweistelligen kHz-Bereich möglich wurden. Die effektive Transportverzögerung 
zwischen Messung und Stellgröße sank dadurch deutlich, was die Stabilitätsreserven des Reglers erhöhte \cite{zander2023}.

Die Balanceregelung realisierte Zander als klassisch ausgelegten PID-Regler, dessen Parameter iterativ auf 
Basis umfangreicher Mess- und Logdaten optimiert wurden. Gegenüber der zuvor genutzten dreistufigen P-Logik 
\cite{karabila2022} steigerten Integral- und D-Anteil Ausregelgüte und Dämpfung. In der Software wurden robuste 
Elemente berücksichtigt, unter anderem Begrenzung und Glättung des Lenkmotor-Solls (Rate- und Amplitudenlimits), 
sorgfältige Filterung der Eingangssignale (Encoder/IMU) zur Rauschunterdrückung bei erhaltener Phasenreserve 
sowie latenzarme Pfade zwischen PRU-Puffern und Reglerkern zur Minimierung von Jitter. Das Loggingsystem 
zeichnete synchronisierte Sensor- und Aktordaten auf und erlaubte eine systematische 
Auswertung der Reglerwirkung \cite{zander2023}.

Neben dem Reglerkern stärkte Zander die Betriebssicherheit: Zwei Notausschalter (separate Lenkungs- bzw. Gesamtabschaltung), 
eine Reißleine zur sofortigen Trennung der Lenkungsversorgung in kritischen Situationen sowie ein Fail-Safe 
für den RC-Empfänger (Stop bei Signalverlust) begrenzen Risiken unkontrollierter Lenkereinschläge und mechanischer Schäden \cite{zander2023}.

In experimentellen Fahrtests blieb das Fahrrad über mehr als 60\,s stabil in Balance, ohne dass die Stützräder 
den Boden berührten; die Zielvorgabe wurde erfüllt und überschritten. Bei konstanter Geschwindigkeit auf ebenem 
Untergrund waren zudem gesteuerte Richtungsänderungen per Fernbedienung möglich, ohne die Balance zu verlieren 
\cite{zander2023}. Damit erzielte Zander den bis dahin längsten autonom balancierten Fahrbetrieb des Prototyps.

Für die vorliegende Arbeit ist diese PID-Basis maßgeblich: Die in \cite{zander2023} entworfene Regelung wurde 
nach \emph{Webots} portiert und bildet die innere Schleife der dualen Reglerarchitektur. Darauf aufbauend wird 
die vision- und trajektorienbasierte Fahreglung in der äußeren Schleife integriert. Diese klare Schichtung 
(Balance innen, Fahrspur außen) ist zentral für die Kopplung von Regelungs- und Computer-Vision-Methoden im Gesamtsystem.

  
\subsubsection{Yasin Giray (2024) – YOLO + MPC-Fahrtrichtung}

Aufbauend auf der in \cite{zander2023} erreichten stabilen Balanceregelung verlagerte Giray (2024) \cite{Ghidaei2024} 
den Schwerpunkt auf die autonome Fahrtrichtungswahl mittels Computer Vision und modellprädiktiver Regelung. Ziel war 
eine monokamera-basierte Pipeline, die in Echtzeit auf eingebetteter Hardware lauffähig ist und aus Bilddaten eine 
Referenztrajektorie ableitet, der das Fahrrad folgen kann, ohne die Balance zu verlieren.

Die Wahrnehmungskomponente setzt auf ein effizientes Segmentierungsmodell aus der YOLOv8-Familie \cite{Ultralytics2023}. 
Nach einer Evaluation alternativer Verfahren (u.\,a. FastSAM \cite{Zhao2023}) fiel die Wahl aufgrund des guten 
Verhältnisses von Genauigkeit und Latenz auf ein Ultralytics-Modell in der Segmentation-Variante. Das Netz wurde 
mit einem angepassten Datensatz trainiert und für die Inferenz mittels TensorRT auf dem NVIDIA Jetson Nano beschleunigt, 
sodass die Bildauswertung in Zyklen von etwa 100\,ms erfolgen konnte. Aus den segmentierten Szenen werden Stützpunkte 
extrahiert und über Spline-Interpolation zu einer glatten Referenzspur verdichtet, die die gewünschte Fahrlinie repräsentiert.

Für die Fahrregelung implementierte Giray einen Model Predictive Controller (MPC) auf Basis eines Einspurmodells 
\cite{Hunziker2019}. Der MPC minimiert über einen endlichen Horizont Abweichungen von der Referenzspur bei gleichzeitig 
moderatem Stellaufwand und berücksichtigt Nebenbedingungen des Systems. In jedem Rechenzyklus wird das Optimierungsproblem 
mit der aktuellen Zustandsschätzung gelöst; angewandt wird jeweils nur die erste Stellgröße (receding horizon), bevor der 
Vorgang mit aktualisierten Messdaten wiederholt wird. Modellierung, Trajektorien-Interpolation und MPC-Optimierung wurden 
zunächst in Python umgesetzt und auf die Jetson-Plattform portiert.

In Simulationen zeigte die Pipeline, dass Fahrbahnkanten und einfache Hindernisse zuverlässig detektiert und Referenzspuren 
generiert werden können; das Fahrrad folgte diesen Spuren stabil, solange die äußere MPC-Schleife die innere Balanceregelung 
nicht überfordert. Sichtbar wurden zugleich Grenzen und Folgearbeiten: Eine explizite Kollisionsvermeidung war noch nicht 
integriert; Tiefe stand nur implizit über Bildgeometrie zur Verfügung. Als Verbesserungen werden daher die Anreicherung um 
Tiefeninformationen (Stereo oder monocular depth) sowie eine kameraseitige Stabilisierung gegen Rollneigung (z.\,B. Gimbal) 
genannt. Ebenso ist mit neueren Segmentierungsansätzen und leistungsfähigeren Edge-Plattformen eine weitere Reduktion der 
Latenz zu erwarten \cite{Zhao2023,Ultralytics2023}.

Für die vorliegende Arbeit ist die Einordnung klar: Die von Zander übernommene 
PID-Balanceregelung bildet die innere Schleife, 
während Girays MPC-basierter Spurhalter als äußere Schleife in \emph{Webots} integriert wird. Diese Schichtung erlaubt es, die 
Kopplung von Balancehaltung und visionbasierter Trajektorienverfolgung systematisch zu untersuchen und in einer realitätsnahen 
Simulation zu kalibrieren, bevor reale Fahrversuche erfolgen.
